\section{Discussion}
\label{sec:discussion}

\subsection{Interpreting the Observed Portfolio Dependence}
\label{sec:mismatch}

Within the evaluated protocol, a fixed diagnostic's measured value changes with the trained models its actions select. Combined with MLP fallback incurs greater descriptive regret than always-graph for the full portfolio, while the post-hoc GCN and GAT restrictions reverse that ordering. The favorable results are narrow: only those two architectures and their pair favor Combined among the 63 subsets, and excluding both Cornell and Wisconsin removes the advantage everywhere. These observations motivate evaluating each diagnostic--portfolio pairing. Their sensitivity to dataset composition leaves performance on independent datasets unresolved.

The full-set Combined action simplifies to a single homophily threshold under the MLP fallback. Roman-empire, Squirrel, Amazon-ratings, and Chameleon account for 95.1\% of its full-portfolio regret: all fall below that threshold, and the rule declines graph on all 40 units. The architecture-selection summary in Table~\ref{tab:winning_architectures} helps locate this mismatch, but selection within the graph portfolio is distinct from beating MLP. The concentration identifies where this complete policy is costly without establishing an architecture mechanism. Separately, 75.9\% of its total regret occurs on the 35 abstaining units that receive MLP fallback. These overlapping groupings locate the losses, not independent causes; Tables~\ref{tab:fallback_sensitivity} and~\ref{tab:preprocessing_fallback} show how their magnitude changes under the specified alternative configurations.

Threshold choice remains consequential. In the current full-portfolio evaluation, the post-hoc held-out calibration reduces regret but still exceeds always-graph (Section~\ref{sec:calibrated_threshold}); even choosing the best one-sided homophily cutoff using all test outcomes cannot improve on that reference in these records (Appendix~\ref{app:threshold_headroom}). Its small 0.256-point headroom limits the possible absolute gain without ruling out gains from richer selectors. The result does not extend to every portfolio: calibrated GCN, GAT, and GPR-GNN have lower descriptive regret than their own always-graph references. Likewise, improvement over always-graph and absolute regret answer different questions: under the frozen rule, GPR-GNN yields lower absolute Combined regret than GCN or GAT while favoring always-graph. The restrictions use different graph/MLP oracles and do not establish a general ordering of architecture categories.

\subsection{Evaluating Decisions and Their Cost}

The frozen comparison evaluates actions on common units with a practical margin and explicit fallback. Its graph-heavy target distribution makes always-graph a strong reference; regret reveals the cost of errors that selection accuracy alone obscures, and coverage exposes abstention. The predeclared incremental-value rule requires lower full-set regret than simple baselines without relying on reduced coverage, with the uncertainty interval excluding zero in the favorable direction for a pre-specified comparison (Section~\ref{sec:protocol}). Combined fails the regret-reduction requirement against always-graph on the recorded point estimates; this does not depend on a Holm significance requirement. The dataset-level tests establish no reference-policy superiority after Holm correction. For the key always-graph comparison, structural zeros and the planned correction make rejection unattainable on this dataset composition, limiting the evidential role of non-rejection. This establishes neither equivalence nor a lack of power for every comparison in the family.

Validation selection achieves the lowest descriptive regret but requires training both model families and is not a guaranteed upper bound. Combined itself uses MLP validation accuracy for coverage, although only homophily determines its full-set actions. A useful selection policy must therefore be assessed against a relevant reference with its decision-time information and training cost stated explicitly.

\subsection{Graph Signal and Diagnostic Value}

The degree-preserving intervention shows useful non-degree edge organization for GCN on Cora, CiteSeer, and PubMed. These three graphs and the intervention architecture differ from the four largest diagnostic-loss settings, so the result supplies auxiliary graph-signal evidence rather than their failure mechanism. Randomization changes several structural attributes together and does not isolate homophily. Appendix~\ref{sec:snr}'s fixed-degree calculation similarly explains interactions among label mixing, degree, and feature strength within a prescribed surrogate; it does not characterize trained graph-model selection.

\subsection{Training Pipelines Define the Comparison}

The learned baseline is part of the estimand. Raw-feature retraining substantially improves MLP accuracy on Roman-empire and improves it on Amazon-ratings, reducing the graph--MLP gaps while preserving the graph target on both. The original gaps consequently combine the effects of model inputs and training choices. Validation-only diagnostics further recover MLP performance through a training-fitted affine reparameterization of normalized inputs, without adding information to those inputs. A common preprocessing transform alone therefore cannot establish an adequately trained feature-only reference.

Training repeatability is a separate concern. Repeated LINKX runs with matching inputs vary substantially in the tested settings. The bounded 22-process audit records default-CUDA variation from identical initial states and random states; its specified deterministic configuration yields identical histories and selected states within the tested repeat groups. Here record reconstruction means recomputing decision scores from saved outcomes; training repeatability means agreement across repeated training executions within a specified configuration. The audit characterizes training repeatability in these tested configurations and does not certify it for the original 770 benchmark records or the 420-record graph diagnostic. The two-dataset raw and centered-scaled diagnostics constrain interpretation; this bounded audit does not replace a complete benchmark rerun or establish that the full benchmark withstands improved baseline training. Appendix~\ref{app:training_diagnostics} preserves these scopes and observed contrasts.

\subsection{Limitations and Next Tests}

The benchmark covers 11 public datasets, one stratified split design, ten seeds, seven architectures, a four-trial grid, and test accuracy. Inference treats datasets as clusters; additional seeds do not supply independent datasets. Thresholds were frozen after inspection of legacy aggregates but before the primary benchmark records. The public snapshot documents file identities and execution provenance, while the earlier private history remains inaccessible.

The graph action receives 24 trials against the MLP's four. Single-architecture restrictions match trial counts while changing architecture availability and search breadth. Equal trial counts do not equalize time or FLOPs; neither these restrictions nor the subset enumeration isolate compute effects or test a comparably expanded MLP search. All reuse the original outcomes. Excluding Chameleon and Squirrel preserves the full-portfolio disadvantage but cannot replace retraining on duplicate-filtered versions.

The next useful evidence is an independent prospective comparison that includes representative published diagnostics and specified candidate portfolios. It should separate selector development from evaluation, plan dataset-level precision and multiplicity, and measure resource consumption alongside trial counts. This would test whether the observed decision effects persist with stronger training baselines and new datasets.
